\documentclass[10pt]{article}
\usepackage[a4paper,margin=1in]{geometry}
\usepackage[T1]{fontenc}
\usepackage{lmodern}
\usepackage[expansion=false]{microtype}
\usepackage{tabularx,array}
\usepackage{enumitem}
\usepackage[hidelinks]{hyperref}
\urlstyle{same}
\usepackage[compact]{titlesec}

% All text stays at 10pt or larger (course minimum); single-spaced, 1-inch margins.
\titleformat*{\section}{\normalsize\bfseries}
\titlespacing*{\section}{0pt}{3pt}{1pt}
\setlength{\parindent}{0pt}
\setlength{\parskip}{1pt}
\setlength{\tabcolsep}{4pt}
\newcolumntype{L}{>{\raggedright\arraybackslash}X}

\begin{document}

% ---------------- Title and team ----------------
\begin{center}
  {\large\bfseries Knowing When the Evidence Is Not Enough:
  Reinforcement Learning for Calibrated Fact Verification\par}
  \vspace{4pt}
  \begin{tabular}{@{}llll@{}}
    \textbf{Name} & \textbf{Matric no.} & \textbf{Email} & \textbf{Role} \\
    Chia Kwang Yang & A0248658L & e0929828@u.nus.edu & AG-1 retrieval \& tools \\
    Yu Jing & A0351724H & jing\_yu@u.nus.edu & AG-2 agent loop \& multi-turn RL \\
    Bryan Wu Jiahe & A0355175Y & wubryan@u.nus.edu & AG-3 evaluation, data \& demo \\
    Chao Chien Hung & A0327301U & e1539132@u.nus.edu & RL-1 verifier training \\
    Michel Yildirim & A0350371M & yldrmmichel@u.nus.edu & RL-2 reward \& calibration \\
    Yao Tianhao & A0274850Y & yao.tianhao@u.nus.edu & RL-3 cost-aware controller \\
  \end{tabular}
\end{center}
\vspace{-8pt}

% ---------------- Body ----------------
\section*{Background}
Platforms are handing fact-checking to AI. Meta ended its US fact-checking programme in April
2025~\cite{meta}, and 447{,}083 posts on X asked Grok for a fact check in six months of 2025, yet the bots
agreed with professional fact-checkers only about half the time~\cite{renault}. A fact-checker that answers
confidently when the evidence is unclear spreads false confidence. Datasets such as FEVER~\cite{fever},
SciFact~\cite{scifact} and AVeriTeC~\cite{averitec} have answer keys, so verdicts can be scored
automatically, which reinforcement learning (RL) needs. RL has been used to make fact-checkers
cite evidence~\cite{verir1} and check claims in steps~\cite{profact}, and to teach honest confidence and
``I don't know'' in question answering~\cite{rlcr,truthrl}. None do this for fact-checking.

\section*{Problem statement}
Given a claim, the model must say whether it is Supported, Refuted or Not Enough Info (NEI), point to the
Wikipedia sentences it relied on, and state how confident it is. Small models that fit on one GPU do this
poorly. On 900 FEVER claims, Qwen3.5-2B is right only 63\% of the time even when shown the correct evidence
(a 180B model: 95\%). It calls almost everything Supported and almost always claims 100\% confidence. Can
RL teach a small model to weigh evidence, say NEI when the evidence falls short, give a confidence that
matches its accuracy, and treat calling a false claim true as the worst mistake?

\section*{Approach}
\textbf{Verifier:} fine-tune the model on labelled examples (SFT), then train it further with RL (GRPO~\cite{grpo}).
The reward pays for correct verdicts, penalises calling false claims true the most, rewards citing the right
sentences, and scores stated confidence against actual accuracy. Varying the NEI reward shows how accuracy
trades off against how often the model abstains, so we can choose how cautious it is. \textbf{Controller:} a
small second model that decides whether to search again, split the claim, verify, abstain or answer, with a
cost for each search. \textbf{Evaluation:} one shared script scores every run for accuracy, confidence
calibration and search cost, and compares RL against the SFT model and simple baselines~\cite{spurious}. We
test for shortcuts with the FEVER symmetric set~\cite{symmetric} and by swapping evidence between claims,
and test on SciFact and AVeriTeC to see if the model generalises.

\section*{Team roles and deadlines}
\begin{tabularx}{\linewidth}{|l|L|L|}
  \hline
  & \textbf{Agentic group (AG-1, AG-2, AG-3)} & \textbf{RL group (RL-1, RL-2, RL-3)} \\ \hline
  Wk 7 & Agent loop, dataset loaders, shortcut tests & Try the 4B model, fix the reward \\ \hline
  Wk 8--9 & Agent baseline, multi-step RL training & RL training, reward experiments \\ \hline
  Wk 10--11 & Tests on new datasets, error analysis, demo & Controller training, cost experiments \\ \hline
  Wk 12 & \multicolumn{2}{l|}{Final report, slides and live demo} \\ \hline
\end{tabularx}

\section*{Risk assessment and preliminary results}
We have already tested the riskiest steps. Fine-tuning took 37 minutes on our 16\,GB GPU (RTX 5070Ti).
\begin{itemize}[leftmargin=*,nosep]
  \item \textbf{The 2B model is too small.} Fine-tuning raised it from 63\% to 82\% with the correct evidence
        shown. In week 7 we try the 4B model (about 1.5 hours as a test) before choosing.
  \item \textbf{RL adds nothing over fine-tuning.} Fine-tuning alone raised correct NEI answers from 14\% to
        72\% (180B model: 35\%). RL targets what fine-tuning cannot fix: the model still always claims 100\%
        confidence, and says NEI on 20\% of false claims even when shown evidence that disproves them.
  \item \textbf{The model takes shortcuts.} Given the wrong claim's evidence, it says NEI on 98\% of true
        claims, so it does read the evidence. However, it still calls 23\% of false claims false from their
        wording alone, so both shortcut tests run on every training run.
  \item \textbf{The model says NEI to everything to play safe.} We report accuracy at every level of caution.
  \item \textbf{It only works on FEVER.} We test on AVeriTeC's real-world claims and report the gap.
  \item \textbf{One GPU for two training jobs.} We share GPU time, and the controller trains on CPU.
\end{itemize}

% ---------------- References (Vancouver style, in order of first citation) ----------------
\clearpage
\begin{thebibliography}{99}\setlength{\itemsep}{1pt}\raggedright
\bibitem{meta} Meta. More speech and fewer mistakes [Internet]. Menlo Park (CA): Meta; 2025 Jan 7
  [cited 2026 Sep 26]. Available from: \url{https://about.fb.com/news/2025/01/meta-more-speech-fewer-mistakes/}
\bibitem{renault} Renault T, Mosleh M, Rand DG. @Grok is this true? LLM-powered fact-checking on social
  media [Preprint]. PsyArXiv; 2025 [cited 2026 Sep 26]. Available from: \url{https://doi.org/10.31234/osf.io/85quw}
\bibitem{fever} Thorne J, Vlachos A, Christodoulopoulos C, Mittal A. FEVER: a large-scale dataset for fact
  extraction and VERification. In: Proceedings of the 2018 Conference of the North American Chapter of the
  Association for Computational Linguistics: Human Language Technologies; 2018 Jun; New Orleans, LA.
  Stroudsburg (PA): Association for Computational Linguistics; 2018. p. 809-19.
\bibitem{scifact} Wadden D, Lin S, Lo K, Wang LL, van Zuylen M, Cohan A, et al. Fact or fiction: verifying
  scientific claims. In: Proceedings of the 2020 Conference on Empirical Methods in Natural Language
  Processing (EMNLP); 2020 Nov; online. Stroudsburg (PA): Association for Computational Linguistics; 2020.
  p. 7534-50.
\bibitem{averitec} Schlichtkrull M, Guo Z, Vlachos A. AVeriTeC: a dataset for real-world claim verification
  with evidence from the web. Adv Neural Inf Process Syst. 2023;36:65128-67.
\bibitem{verir1} He Q, Qian C, Chen X, He B, Fung YR, Ji H. Veri-R1: toward precise and faithful claim
  verification via online reinforcement learning [Preprint]. arXiv:2510.01932. 2025 [cited 2026 Sep 26].
  Available from: \url{https://arxiv.org/abs/2510.01932}
\bibitem{profact} Yang R, He S, Zhu S, Yu C. From verdict to process: agentic reinforcement learning for
  multi-stage fact verification [Preprint]. arXiv:2606.13262. 2026 [cited 2026 Sep 26].
  Available from: \url{https://arxiv.org/abs/2606.13262}
\bibitem{rlcr} Damani M, Puri I, Slocum S, Shenfeld I, Choshen L, Kim Y, et al. Beyond binary rewards:
  training LMs to reason about their uncertainty [Preprint]. arXiv:2507.16806. 2025 [cited 2026 Sep 26].
  Available from: \url{https://arxiv.org/abs/2507.16806}
\bibitem{truthrl} Wei Z, Yang X, Sun K, Wang J, Shao R, Chen J, et al. TruthRL: incentivizing truthful LLMs
  via reinforcement learning [Preprint]. arXiv:2509.25760. 2025 [cited 2026 Sep 26].
  Available from: \url{https://arxiv.org/abs/2509.25760}
\bibitem{grpo} Shao Z, Wang P, Zhu Q, Xu R, Song J, Bi X, et al. DeepSeekMath: pushing the limits of
  mathematical reasoning in open language models [Preprint]. arXiv:2402.03300. 2024 [cited 2026 Sep 26].
  Available from: \url{https://arxiv.org/abs/2402.03300}
\bibitem{spurious} Shao R, Li SS, Xin R, Geng S, Wang Y, Oh S, et al. Spurious rewards: rethinking training
  signals in RLVR [Preprint]. arXiv:2506.10947. 2025 [cited 2026 Sep 26].
  Available from: \url{https://arxiv.org/abs/2506.10947}
\bibitem{symmetric} Schuster T, Shah D, Yeo YJS, Filizzola Ortiz DR, Santus E, Barzilay R. Towards debiasing
  fact verification models. In: Proceedings of the 2019 Conference on Empirical Methods in Natural Language
  Processing and the 9th International Joint Conference on Natural Language Processing (EMNLP-IJCNLP);
  2019 Nov; Hong Kong, China. Stroudsburg (PA): Association for Computational Linguistics; 2019. p. 3419-25.
\end{thebibliography}

\end{document}
